\section{条件性实施研究计划}

本章给出条件性模型仓位的建立顺序，不构成交易执行授权。未补齐原始证据、未取得实施日产品和市场数据并完成独立复核前，不得据此提交交易委托或确认价格、数量、通道或适当性。没有完整执行日资格包时，所有ETF资格开关均为$q=0$；各笔目标、条件及失败后的资金位置仍可复核。

\begin{exhibitbox}[六个月模型实施动作]
\footnotesize
\begin{tabularx}{\exhibitboxwidth}{L{2.6cm}P P}
\toprule
\textbf{阶段} & \textbf{模型金额} & \textbf{应完成的检查与失败处理} \\
\midrule
启动防守层 & CNY40,000受保障流动性、CNY10,000交易现金和CNY70,000的\texttt{511010} & 验证资金用途、存款保险边界、账户和产品资料；\texttt{511010}还须通过当前价格/IOPV、折溢价、价差和流动性检查。未通过的ETF金额留在受保障流动性。 \\
第1--5个月 & 每月CNY5,000的\texttt{510880}、CNY5,000的合格\texttt{512890}和CNY3,333的\texttt{518880} & 每笔记录同一市场时段的来源、时间戳、价格/IOPV、价差和流动性；\texttt{512890}另需合并持仓穿透。任一失败只延后该笔，并留在受保障流动性。 \\
第6个月 & CNY5,000的\texttt{510880}、CNY5,000的合格\texttt{512890}和CNY3,335的\texttt{518880} & 以最小交易单位向下取整；剩余与未通过金额留在受保障流动性，交易现金仍不得超过CNY10,000。 \\
REITs与直接股票 & CNY0 & 未通过独立激活门槛时保持0，不以市场下跌、等待分配或ETF延后绕过资格开关。 \\
\bottomrule
\end{tabularx}
\end{exhibitbox}

金额为风险预算的算术映射，唯一后备位置是受保障流动性：核验失败不会增加\texttt{511010}、替换ETF或令交易现金超过CNY10,000。延后额度只可在下一计划分批或再平衡日重取完整资格包后实施，不补单。

\section{实施日前的资格检查}

{\footnotesize
\begin{enumerate}
  \item 确认固定CNY200,000确属五年以上可承受波动的本金，且没有改变为刚性用途。
  \item 对每个候选ETF保存当前产品文件、费用、来源、时间戳、交易状态、价格/IOPV、折溢价、最佳买卖价、价差、同日成交额或可执行深度、账户资格、最小交易单位和现金可行性。
  \item ETF仅在绝对折溢价不高于1.00\%、买卖价差不高于中间价0.50\%、单笔不高于记录同日成交额或可执行深度10\%时$q=1$；缺字段或未通过即$q=0$。
  \item 对两只红利权益ETF做带日期的合并持仓穿透，确认行业不超过总组合25\%，前十大集团不超过权益桶35\%；\texttt{512890}未通过时保持0。
  \item 核验利率债久期与用途匹配且不把债券/货币ETF写成保本存款；黄金只可为非杠杆现货ETF；公募REITs须有至少三只独立项目通过现金流、债务、估值溢价、流动性和最小交易单位审查，才可从利率债转入最多5\%。
\end{enumerate}
\par}

\section{再平衡与记录}

{\footnotesize
\begin{itemize}
  \item \textbf{季度}：记录实际到账的税后、扣费后现金$D_0$、资产权重、产品事件、资格开关、外部现金流和实际$DD_t$；$D_0$须有到账日、来源、税费、净额和凭证。
  \item \textbf{半年}：优先用新增资金和实际现金分配恢复权重；受保障流动性、利率债或总权益偏离5个百分点、交易现金超过CNY10,000或分批期内低于目标、黄金偏离2个百分点时复核。
  \item \textbf{实际回撤与年度}：已记录$DD_t\leq-12\%$时停止自动新增风险暴露，$DD_t\leq-18\%$时先复核流动性、杠杆、产品异常和资金用途；情景损失不替代$DD_t$。年度重读产品文件和税务规则，更新现金、穿透和观察池；未独立研究时直接股票仍为0\%。
\end{itemize}
\par}

\qualitynote{本文件为条件性研究草案。实施日资格包、个人账户适当性、税务和任何交易决定须由投资者及其受托专业人士独立确认。}
